\DSource{0132}{43}
\hypertarget{AB01-PDF0132-H01}{}
\DLine{AB01-PDF0132-body-L001}{}{}{\CenterLine{CAPUT XXVIII.}}
\vspace{13.8pt}
\DLine{AB01-PDF0132-body-L002}{5}{}{\CenterLine{De inaequalitate motus [sive anomalia] Solis et de eo quod}}
\DLine{AB01-PDF0132-body-L003}{}{}{\CenterLine{ab eius apogei loco cum illa apparet (1).}}
\vspace{13.8pt}
\hypertarget{AB01-PDF0132-P01}{}
\DLine{AB01-PDF0132-body-L004}{}{}{\Indent Dicit [auctor]: Dissertatione de tempore anni et de motu medio Solis absoluta, ad expli-}
\DLine{AB01-PDF0132-body-L005}{}{}{candum pergamus quae et quanta sit inaequalitas in motu Solis, et qualis appareat [cum Sol}
\DLine{AB01-PDF0132-body-L006}{10}{}{motu suo medio recesserit] a loco sui apogei in ecliptica.}
\hypertarget{AB01-PDF0132-P02}{}
\DLine{AB01-PDF0132-body-L007}{}{}{\Indent Rationem qua Ptolemaeus in libro suo (2) usus est sequentes, nos quo modo Sol quadran-}
\DLine{AB01-PDF0132-body-L008}{}{}{tes eclipticae percurrat, per plurimas diligentissimas observationes multis subsequentibus annis}
\DLine{AB01-PDF0132-body-L009}{}{}{factas studuimus, invenimusque Solem eclipticam a puncto aequinoctii autumni ad aequinoctii}
\DLine{AB01-PDF0132-body-L010}{}{}{vernalis 178 diebus, 14 horis $\qfrac{1}{2}$ circiter percurrere; a puncto autem aequinoctii vernalis ad}
\DLine{AB01-PDF0132-body-L011}{15}{p. 65.}{aequinoctii autumni tempus eo longius occupare, quod ex nostris diligentissimis observationibus}
\DLine{AB01-PDF0132-body-L012}{}{}{patuit 186 dierum, 14 horarum aequinoctialium et $\qfrac{1}{4}+\qfrac{1}{2}$ horae circiter esse. Et ex iis}
\DLine{AB01-PDF0132-body-L013}{}{}{quae memoravimus (3) apparuit eius apogeum in hac dimidia parte haberi.}
\hypertarget{AB01-PDF0132-P03}{}
\DLine{AB01-PDF0132-body-L014}{}{}{\Indent Deinde observavimus et invenimus Solem spatium ab initio Arietis ad initium Cancri,}
\DLine{AB01-PDF0132-body-L015}{}{}{idest inter punctum aequinoctii vernalis et punctum solstitii aestivi, 93 diebus et 14 fere horis}
\DLine{AB01-PDF0132-body-L016}{20}{}{aequinoctialibus percurrere ({\itshape a}), quae parum ad deminutionem inclinant; unde manifestum est}
\DLine{AB01-PDF0132-body-L017}{}{}{cursui Solis a puncto aequinoctii vernalis ad solstitii aestivi tempus longius necessarium esse}
\DLine{AB01-PDF0132-body-L018}{}{}{quam a puncto solstitii aestivi ad aequinoctii autumni. Scimus igitur punctum apogei et cen-}
\DLine{AB01-PDF0132-body-L019}{}{}{trum excentrici, super quem punctum apogei simul ac super eclipticam incidit, ea quarta parte}
\DLine{AB01-PDF0132-body-L020}{}{}{contineri, quae tardioris temporis est quam reliquus quadrans. Motum medium Solis in iis}
\DLine{AB01-PDF0132-body-L021}{25}{}{186 diebus, 14 horis et $\qfrac{1}{2}+\qfrac{1}{4}$ horae invenimus $183^\circ56'12''$ esse, et in iis 93 diebus et}
\DLine{AB01-PDF0132-body-L022}{}{}{14 horis $92^\circ14'10''$ fere (4).}
\hypertarget{AB01-PDF0132-P04}{}
\DLine{AB01-PDF0132-body-L023}{}{}{\Indent Quae cum ita sint, circulum eclipticae ABCD circa centrum E describamus. Duae dia-}
\DLine{AB01-PDF0132-body-L024}{}{}{metri AC, BD se invicem ad angulos rectos abscindant, sitque A punctum aequinoctii vernalis;}
\DLine{AB01-PDF0132-body-L025}{}{}{erit B punctum solstitii aestivi, C aequinoctii autumni, D solstitii hiemalis. In quadrante AB,}
\DLine{AB01-PDF0132-body-L026}{30}{}{ob ea quae supra diximus, punctum F centrum ponamus, et, intra primum circulum, circu-}
\DLine{AB01-PDF0132-body-L027}{}{}{lum KLMN sphaerae excentricae Solis describamus, cuius diametri KM, LN se invicem in}
\DLine{AB01-PDF0132-body-L028}{}{}{centro F ad angulos rectos abscindant. Punctum lineis BD, KM commune littera S signemus;}
\DLine{AB01-PDF0132-body-L029}{}{}{punctum quo diametrus AC circulum KLMN versus A secat, littera P; punctum quo diame-}
\DLine{AB01-PDF0132-body-L030}{}{}{trus BD circulum KLMN versus B abscindit, littera R. In arcu PK, perpendicularem a P ad}
\DLine{AB01-PDF0132-body-L031}{35}{}{punctum Q diametri KM ducamus; item perpendicularem RG, et lineam EF quae per duo}
\DLine{AB01-PDF0132-body-L032}{}{}{centra transit, eamque usque ad punctum H eclipticae ABCD protrahamus, littera T signantes}
\DLine{AB01-PDF0132-body-L033}{}{}{locum quo ea circulum KLMN secat.}
\hypertarget{AB01-PDF0132-P05}{}
\DLine{AB01-PDF0132-body-L034}{}{}{\Indent Manifestum est arcum AB esse $90^\circ$, ut arcum KL excentrici; punctum P excentrici esse}
\DLine{AB01-PDF0132-body-L035}{}{p. 66.}{initium Arietis; R locum initii Cancri; denique arcum PKLRMY excentrici eam excentrici}
\par\vspace{6pt}\hrule height.25pt\vspace{5pt}
\noindent\begin{minipage}[t]{.48\textwidth}\vspace{0pt}\fontsize{9}{11.5}\selectfont\setlength{\GutterOffset}{0pt}
\hypertarget{AB01-PDF0132-N01}{}
\DLine{AB01-PDF0132-notes_left-L001}{}{}{\Indent (1) I. e. de supputanda anomalia ad singulos}
\DLine{AB01-PDF0132-notes_left-L002}{}{}{gradus medii motus Solis ab apogeo.}
\hypertarget{AB01-PDF0132-N02}{}
\DLine{AB01-PDF0132-notes_left-L003}{}{}{\Indent (2) {\itshape Almag.} III, 4 (ed. Halma t. I, p. 184--185.}
\hypertarget{AB01-PDF0132-N03}{}
\DLine{AB01-PDF0132-notes_left-L004}{}{}{\Indent (3) Sed antea, ut me \Name{Schiaparelli} monet,}
\DLine{AB01-PDF0132-notes_left-L005}{}{}{nihil de his rebus dixit. Forte aliquid excidit ubi}
\DLine{AB01-PDF0132-notes_left-L006}{}{}{demonstrabatur apogeum reperiendum esse in parte}
\DLine{AB01-PDF0132-notes_left-L007}{}{}{cui maius temporis intervallum respondet.}
\end{minipage}
\hfill\begin{minipage}[t]{.48\textwidth}\vspace{0pt}\fontsize{9}{11.5}\selectfont\setlength{\GutterOffset}{0pt}
\hypertarget{AB01-PDF0132-N04}{}
\DLine{AB01-PDF0132-notes_right-L001}{}{}{\Indent (4) Ita etiam in calculis infra, qua re corri-}
\DLine{AB01-PDF0132-notes_right-L002}{}{}{gere non audeo; e tabulis $92^\circ14'26''$ deducuntur.}
\DLine{AB01-PDF0132-notes_right-L003}{}{}{Sed antea dixerat « 93 diebus et 14 horis aequino-}
\DLine{AB01-PDF0132-notes_right-L004}{}{}{« ctialibus, {\itshape quae parum ad deminutionem incli-}}
\DLine{AB01-PDF0132-notes_right-L005}{}{}{« {\itshape nant} »; forte huic « parum minus » respondent}
\DLine{AB01-PDF0132-notes_right-L006}{}{}{hae 16 secundae, quae in auctoris calculo deficere}
\DLine{AB01-PDF0132-notes_right-L007}{}{}{videntur.}
\end{minipage}
